\textbf{Abstract} Collective effects in modern particle accelerators generate nonlinear correlations in the six-dimensional beam phase space and require computationally expensive multiparticle simulations. This work investigates exploratory surrogate models for beam-distribution evolution in the High-Energy Booster of the Future Circular Collider (FCC-ee), with the objective of accelerating numerical studies while retaining physically relevant distributional information. Two complementary modelling strategies are developed. Reduced longitudinal dynamics are learned directly as mappings between density functions using Fourier Neural Operators. For the six-dimensional problem, beam states are represented by the fifteen pairwise two-dimensional phase-space marginals, supplemented by six centroids and six coordinate-wise RMS spreads. A conditional variational autoencoder compresses these projected distributions into a probabilistic latent representation, whose full-lattice input-to-output evolution is modelled using a causal Transformer, a local-window Transformer, or a Graph Neural Operator. Only the final downstream state is supervised and validated; the internal element-indexed latent states are not claimed to be physically validated intermediate beam states. The comparison of the latent propagators shows that restricting the support of information transfer provides a beneficial inductive bias for the investigated endpoint-transport task. Relative to global causal attention, local-window attention reduces distributional reconstruction errors and systematic artefacts. Among the selected checkpoints, the Graph Neural Operator has the lowest observed relative histogram error, decreasing it from \(0.681\) for the global Transformer to \(0.511 ,\) while requiring approximately \(2.34\) ms per sample under the reported neural-model evaluation conditions, about eight times faster than the global Transformer. These values come from selected single-seed runs and are not accompanied by seed-to-seed uncertainty estimates. The conclusions remain restricted to the sampled simulation domain, the fixed full-lattice endpoint map, and the fifteen-marginal representation, which does not uniquely determine the complete six-dimensional joint distribution. Conditional latent sampling produces structured variability but is strongly under-dispersed relative to the observed prediction error and must not be interpreted as calibrated predictive coverage. Within these limitations, the study identifies structured causal support, explicit physical-statistics prediction, and graph-based latent propagation as promising components for future accelerator surrogate models. Preliminary results from this work were published in the JACoW proceedings of the International Particle Accelerator Conference (IPAC) \cite{ghribi2026updates}. \textbf{Keywords} \emph{Machine} \emph{learning,} \emph{accelerator} \emph{physics,} \emph{FCC-ee,} \emph{collective} \emph{effects,} \emph{beam} \emph{dynamics,} \emph{neural} \emph{oper-} \emph{ators,} \emph{Fourier} \emph{Neural} \emph{Operator,} \emph{conditional} \emph{variational} \emph{autoencoder,} \emph{Graph} \emph{Neural} \emph{Operator,} \emph{Transformer,} \emph{surrogate} \emph{modelling.}
